\section*{الفصل \ref{ch:g6:rocks-and-minerals} --- الصخور والمعادن}

\begin{solution}{exo:g6:rocks-and-minerals:1}
الغرانيت: ناري. الحجر الرملي: رسوبي. الرخام: متحول. البازلت:
ناري. الحجر الجيري: رسوبي. الأردواز: متحول.
\end{solution}

\begin{solution}{exo:g6:rocks-and-minerals:2}
الكوارتز \omterm{def:g6:rocks-and-minerals:mineral}{معدن}؛ والغرانيت \omterm{def:g3:stones-pebbles-sand:rock}{صخر}. \omterm{def:g6:rocks-and-minerals:mineral}{فالمعدن} مادة طبيعية واحدة ذراتها مرتبة
ترتيبًا منتظمًا؛ أما \omterm{def:g3:stones-pebbles-sand:rock}{الصخر} فيتكون من \omterm{def:g6:rocks-and-minerals:mineral}{معدن} واحد أو من عدة \omterm{def:g6:rocks-and-minerals:mineral}{معادن} ملتصقة
بعضها ببعض (الغرانيت: فلسبار وكوارتز وميكا).
\end{solution}

\begin{solution}{exo:g6:rocks-and-minerals:3}
التلك (1)، الكالسيت (3)، الكوارتز (7)، الماس (10). والماس، وهو الأصلد،
يخدشها كلها.
\end{solution}

\begin{solution}{exo:g6:rocks-and-minerals:4}
«يصير \omterm{def:g3:stones-pebbles-sand:rock}{الصخر} المنصهر \omterm{def:g3:stones-pebbles-sand:rock}{صخرًا} ناريًا عن طريق التبرد. ويصير الحجر الجيري
رخامًا عن طريق الحرارة والضغط. ويتحول الغرانيت إلى حبيبات \omterm{def:g3:stones-pebbles-sand:sand}{رمل}
عن طريق \omterm{def:g5:water-shapes-the-land:weathering}{التجوية}.»
\end{solution}

\begin{solution}{exo:g6:rocks-and-minerals:5}
(أ) البازلت. (ب) الرخام. (ج) الحجر الرملي.
\end{solution}

\begin{solution}{exo:g6:rocks-and-minerals:6}
\omterm{def:g6:rocks-and-minerals:hardness}{صلادته} أكبر من \omterm{def:g6:rocks-and-minerals:hardness}{صلادة} الظفر (نحو 2.5) وأقل من \omterm{def:g6:rocks-and-minerals:hardness}{صلادة} قطعة النقود
النحاسية (نحو 3.5). قد يكون الكالسيت (3).
\end{solution}

\begin{solution}{exo:g6:rocks-and-minerals:7}
الكوارتز: $21 \div 60 = 0.35$، أي 35~\%. الفلسبار: $33 \div 60 = 0.55$،
أي 55~\%. الميكا: $6 \div 60 = 0.10$، أي 10~\%. التحقق:
$35 + 55 + 10 = 100$.
\end{solution}

\begin{solution}{exo:g6:rocks-and-minerals:8}
يبرد الغرانيت ببطء شديد، في أعماق الأرض: فتجد حبيباته وقتًا طويلًا لتنمو،
فتصير كبيرة. أما البازلت فيأتي من \omterm{def:g4:volcanoes-earthquakes-first-look:eruption}{الحمم}، التي تبرد بسرعة على السطح: فلا تجد
حبيباته وقتًا لتنمو، وتبقى أصغر من أن تُرى.
\end{solution}

\begin{solution}{exo:g6:rocks-and-minerals:9}
(أ) الحجر الجيري رسوبي؛ والرخام متحول. (ب) في أعماق الأرض، حوّلته الحرارة
والضغط دون أن يصهراه: فنما الكالسيت الذي فيه بلوراتٍ لامعة متداخلة.
(ج) كلاهما مكوَّن من \omterm{def:g6:rocks-and-minerals:mineral}{المعدن} نفسه، الكالسيت، والكالسيت هو الذي يُحدث
فورانًا تحت الحمض.
\end{solution}

\begin{solution}{exo:g6:rocks-and-minerals:10}
(أ) نعم: الكوارتز (7) أصلد من الزجاج (نحو 5.5). (ب) لا: الفلوريت (4)
أطرى من الزجاج. (ج) نعم: المسمار الفولاذي (نحو 6.5) أصلد من
الفلسبار (6).
\end{solution}

\begin{solution}{exo:g6:rocks-and-minerals:11}
مثلًا: (1) تفتت \omterm{def:g5:water-shapes-the-land:weathering}{التجوية} غرانيت \omterm{def:g1:landscapes-around-us:landform}{الجبل} (وهو \omterm{def:g6:rocks-and-minerals:rock-families}{صخر ناري})، فتتحرر حبيبات
الكوارتز التي فيه؛ (2) يحملها \omterm{def:g1:landscapes-around-us:water}{نهر} بعيدًا: إنها \omterm{def:g5:water-shapes-the-land:erosion}{التعرية}؛ (3) تترسب
الحبيبات في طبقات في قاع \omterm{def:g2:water-on-earth:ocean}{البحر}: إنه \omterm{def:g5:water-shapes-the-land:deposition}{الترسيب}؛ (4) تُدفن تحت طبقات أخرى،
فتنضغط وتلتحم حتى تصير حجرًا رمليًا: إنه التصلب (\omterm{def:g3:stones-pebbles-sand:rock}{والصخر} الآن رسوبي)؛
(5) يُحمل الحجر الرملي إلى أعماق الأرض، فتحوّله الحرارة والضغط
إلى \omterm{def:g6:rocks-and-minerals:rock-families}{صخر متحول}.
\end{solution}

\begin{solution}{exo:g6:rocks-and-minerals:12}
(أ) لا يخدشه الفولاذ (نحو 6.5) ويستطيع أن يخدش الزجاج: فهو أصلد
من 6.5. ويخدشه التوباز (8): فهو أطرى من 8. \omterm{def:g6:rocks-and-minerals:hardness}{صلادته} 7: إنه
الكوارتز. (ب) الكوارتز أصلد من كل ما يحتك به تقريبًا
(الحبيبات الأخرى، \omterm{def:g3:stones-pebbles-sand:pebble}{والحصى}): فلا يكاد يبلى. أما الكالسيت (3) والميكا
(من 2 إلى 3) فأطرى منه بكثير: فالاحتكاك يطحنهما قطعًا دقيقة جدًا
يحملها الماء بعيدًا.
\end{solution}

\begin{solution}{pb:g6:rocks-and-minerals:1}
\textbf{1.} الغرانيت، وهو \omterm{def:g6:rocks-and-minerals:rock-families}{صخر ناري}؛ ومعادنه الفلسبار (وردي أو
أبيض)، والكوارتز (رمادي، زجاجي)، والميكا (سوداء، لامعة).

\textbf{2.} الحجر الجيري، وهو \omterm{def:g6:rocks-and-minerals:rock-families}{صخر رسوبي}. كانت الأصداف لحيوانات
عاشت في \omterm{def:g2:water-on-earth:ocean}{البحر}؛ وقد ترسبت بقاياها ثم تصلبت فصارت جزءًا
من \omterm{def:g3:stones-pebbles-sand:rock}{الصخر}.

\textbf{3.} الأردواز، وهو \omterm{def:g6:rocks-and-minerals:rock-families}{صخر متحول}، تكوّن من السجّيل.

\textbf{4.} الرخام، وهو \omterm{def:g6:rocks-and-minerals:rock-families}{صخر متحول}، تكوّن من الحجر الجيري.

\textbf{5.} الكوارتز: لا يخدشه الفولاذ (نحو 6.5)، وهو يخدش الزجاج
(نحو 5.5)؛ \omterm{def:g6:rocks-and-minerals:hardness}{صلادته} 7.

\textbf{6.} يخدشها الفولاذ (نحو 6.5) ولا يخدشها الزجاج (نحو 5.5):
\omterm{def:g6:rocks-and-minerals:hardness}{صلادتها} بين 5.5 و6.5. إنه الفلسبار (الأورثوكلاز، 6).

\textbf{7.} الميكا (\omterm{def:g6:rocks-and-minerals:hardness}{صلادتها} من 2 إلى 3).

\textbf{8.} نعم: البلاط من الحجر الجيري، المكوَّن من الكالسيت (3)، وهو أطرى من
قطعة النقود النحاسية (نحو 3.5).

\textbf{9.} \omterm{def:g5:water-shapes-the-land:weathering}{التجوية}: \omterm{def:g4:clouds-rain-snow:precipitation}{المطر} يُبلي سطح الحجر ببطء
شديد.

\textbf{10.} تتلف \omterm{def:g5:water-shapes-the-land:weathering}{التجوية} \omterm{def:g1:landscapes-around-us:landform}{الجبل}؛ فتحمل \omterm{def:g1:landscapes-around-us:water}{الأنهار} حبيبات الكوارتز التي فيه وترسّبها في طبقات، في قاع \omterm{def:g1:landscapes-around-us:water}{بحيرة} أو \omterm{def:g2:water-on-earth:ocean}{بحر}؛ وحين تُدفن
وتتصلب، يمكن أن تصير حجرًا رمليًا، وهو \omterm{def:g6:rocks-and-minerals:rock-families}{صخر رسوبي}.

\textbf{11.} العائلات الثلاث كلها موجودة: النارية (غرانيت
النافورة)، والرسوبية (الحجر الجيري للبلاط)، والمتحولة (أردواز
السقف ورخام التمثال).

\textbf{12.} أصلد \omterm{def:g6:rocks-and-minerals:mineral}{معدن} هو كوارتز الغرانيت: \omterm{def:g6:rocks-and-minerals:hardness}{صلادته} على \omterm{def:g6:rocks-and-minerals:hardness}{سلم موس}
\textbf{7}.
\end{solution}
